% !TEX program = pdflatex
\documentclass[a4paper,twocolumn,10pt]{article}

\usepackage[top=30mm,bottom=25mm,left=20mm,right=20mm]{geometry}
\setlength{\columnsep}{7mm}
\usepackage{CJKutf8}
% CJKutf8's min family has no bold shape under pdfLaTeX; reuse regular min.
\DeclareFontShape{C70}{min}{b}{n}{<->ssub*min/m/n}{}
\DeclareFontShape{C70}{min}{bx}{n}{<->ssub*min/m/n}{}
\usepackage{graphicx}
\usepackage{booktabs}
\usepackage{array}
\usepackage{url}
\usepackage{titlesec}
\usepackage{setspace}
\titlespacing*{\section}{0pt}{1.2ex plus .3ex minus .2ex}{0.7ex}
\titlespacing*{\subsection}{0pt}{0.9ex plus .2ex minus .2ex}{0.5ex}
\titleformat{\section}{\large\bfseries\sffamily}{\thesection}{0.7em}{}
\titleformat{\subsection}{\normalsize\bfseries\sffamily}{\thesubsection}{0.6em}{}
\setstretch{1.08}
\pagestyle{empty}

\newcommand{\code}[1]{{\small\ttfamily\bfseries #1}}

\begin{document}
\begin{CJK}{UTF8}{min}

\title{\vspace{-8mm}SOC調査起点からの証跡付き行動列復元に向けた\\
エージェント型ログ探索手法の評価\\[1mm]
{\large Evaluation of Agent-based Log Investigation for Evidence-backed Behavior Reconstruction from SOC Investigation Clues}}
\author{著者名$^{1}$ \quad 共著者名$^{2}$\\
{\small $^{1}$所属名 \quad $^{2}$所属名}}
\date{}
\maketitle
\thispagestyle{empty}

\section{はじめに}

SOC（Security Operation Center）におけるエンドポイント調査では，検知アラートの名称や短い検知理由だけでは，端末上で実際にどのような行動が発生したかを十分に説明できない場合がある．例えば，Windowsの永続化手法であるRegistry Run Keyへの登録では，調査者が確認したい対象は「永続化の疑い」というラベルそのものではなく，どのプロセスが起動し，どのコマンドラインで，どのレジストリ値を照会または変更したかである．MITRE ATT\&CKではRegistry Run Keys / Startup FolderがBoot or Logon Autostart Executionの一手法として整理されており，Run Keyに登録されたプログラムはユーザログオン時に実行され得る\cite{mitre_t1547_001}．したがって，アラート後の調査では，最終的な悪性・良性判定だけでなく，観測ログに基づく行動列の再構成が重要となる．

このような行動列の復元では，プロセス生成，親子関係，コマンドライン，レジストリ操作，関連するEDR telemetryなど，複数のログソースを横断する必要がある．一方で，SOCで得られる初期手がかりは必ずしも十分ではない．調査者は，host，process，timestampのような少数の起点，またはEDRアラート要約の一部から探索を開始することが多い．LLMを用いたエージェント型調査は，このような少数の手がかりから検索方針を立て，ログを反復的に探索できる可能性を持つ．しかし，モデルがアラート名を言い換えるだけで，ログ上の証跡に基づかない行動を出力する危険もある．

本研究の目的は，少数のSOC調査起点からWindowsエンドポイントログを探索し，調査者が検証可能な証跡付き行動列を復元できるかを評価することである．本稿では，行動列を，主体，操作，対象，コマンドライン，根拠ログを含むステップ列として扱い，エージェント型パイプラインの出力を構造化された行動ステップ列および短い行動系列として評価する．対象事例はDiscord Run Key persistenceであり，現段階では正式な複数事例実験の前段階として，単一事例に対するプレ実験結果を報告する．

本稿の貢献は次の3点である．第一に，SOC調査起点から行動列を復元するエージェント型ログ探索手法の評価枠組みを設計し，入力情報の強さが異なる複数条件で予備的に確認する．第二に，評価単位を自然文の類似度ではなく，行動ステップを構成する必須項目に分解し，復元率，行動順序，証跡同定，過剰主張を測定する．第三に，EDRアラート要約が初期入力として与えられる場合，DB内の探索対象としてのみ存在する場合，およびDBから除外される場合を比較し，単一事例において復元に寄与する情報源を分析する．

\section{関連研究}

\subsection{エンドポイントログを用いた調査支援}

エンドポイント調査では，Windows Securityログ，Sysmon，EDR telemetry，DNS，ブラウザ履歴など，異なる粒度のログを統合して分析する．Sysmonはプロセス生成時のコマンドライン，親プロセス情報，ProcessGUID，ネットワーク接続，レジストリ変更などを記録できる\cite{sysmon}．これらのログは，攻撃行動の観測証跡を得る上で有用である一方，ログソースごとに時刻，識別子，記録項目が異なるため，人手で一貫した行動列にまとめるには時間を要する．

既存のアラート調査では，アラート名やルール説明が調査の入口として用いられる．しかし，アラート要約は調査対象を絞り込むための情報であり，それ自体が観測された行動を完全に表すわけではない．特に，アラート理由をそのまま「行動」として出力すると，ログ証跡に基づく説明ではなく，検知ロジックの再記述になってしまう．本研究では，アラート情報を探索の入口または証跡の一部として利用しつつ，最終出力は観測ログで確認できる行動に限定する．

\subsection{LLMエージェントによるログ探索}

LLMを用いたエージェント型手法は，自然言語の手がかりから調査方針を生成し，検索結果に応じて次の質問を変更できる点に特徴がある．ReActは，推論と外部環境への行動を交互に行う枠組みを示し，知識源や環境と相互作用しながら回答を構成する有効性を報告している\cite{react}．ログ調査でも，初期手がかりから関連時刻範囲やプロセス木を探索し，得られた行をもとに追加のSQL問い合わせを行うような反復が必要になる．この処理は，単発の分類や要約とは異なり，仮説生成，検索，証跡確認，統合の段階を持つ．

一方，LLMエージェントの出力には，近傍の無関係な行動の混入や，証跡で支えられない主張が含まれる可能性がある．そのため，評価では「正解要素を含んだか」だけでなく，「過剰な主張をどれだけ含んだか」，「行動順序を正しく構成したか」，「根拠ログを対応付けられたか」を分けて見る必要がある．本研究では，この点を行動主張単位で評価する．

\subsection{ログ異常検知とデータセット}

ログを用いたセキュリティ分析では，正常系列からの逸脱を検出する研究が多く行われてきた．DeepLogは，システムログを系列として扱い，LSTMにより異常検知と診断を行う手法を提案している\cite{deeplog}．これらの研究は異常の検出に重点を置く一方，本研究は検知後のSOC調査において，観測された行動列を証跡付きで復元する点に焦点を置く．また，本研究で扱うログはATLASv2由来のWindowsホストログである．ATLASは，audit logから攻撃ストーリを復元するために因果依存グラフと系列学習を組み合わせた攻撃調査支援手法である\cite{atlas}．ATLASv2は，ATLASの攻撃シナリオを再現しつつ，通常業務活動とSysmonおよびCarbon Black Cloud由来のログを拡充したデータセットとして報告されている\cite{atlasv2}．本研究では，この種のデータセットを，異常検知そのものではなく，検知後調査の行動復元評価に用いる．

\section{提案手法}

\subsection{パイプライン構成}

図\ref{fig:pipeline}に，本研究で用いるエージェント型ログ探索パイプラインの概要を示す．入力は，SOC clue，CBC alert，またはprocess-time clueである．Chief agentは，初期手がかりから何を確認すべきかという調査方針を生成する．Investigation agentは，その方針をQAAgent / SQL expertへ渡す具体的な質問に変換する．QAAgent / SQL expertは，SQLite化されたエンドポイントログを検索し，timestamp，source stream，PID/PPID，process tree，command line，registry path，source row idなどの観測値を返す．Final synthesisは，得られた観測結果を統合し，主行動列と近傍行動を分離して行動ステップ列と短縮行動系列を生成する．

\begin{figure*}[t]
\centering
\fbox{\begin{minipage}{0.92\linewidth}
\small
\centering
SOC clue / CBC alert / process-time clue\\
$\downarrow$\\
Chief agent: 調査方針を生成\\
$\downarrow$\\
Investigation agent: 確認すべき事実を質問へ変換\\
$\downarrow$\\
QAAgent / SQL expert: SQLite endpoint logsから観測行を取得\\
$\downarrow$\\
Final synthesis: 主行動列と近傍行動を分離\\
$\downarrow$\\
行動ステップ列 + 短縮行動系列 + limitations\\
$\downarrow$\\
行動主張単位の採点
\end{minipage}}
\caption{SOC調査起点から証跡付き行動列を復元するパイプライン}
\label{fig:pipeline}
\end{figure*}

本手法の特徴は，仮説生成と観測ログ検索を分離する点にある．Chief agentはSQLを直接書かず，アラート名やコマンドラインを仮定しない．Investigation agentは，前回までの検索で得られた観測値から次の確認事項を決める．QAAgent / SQL expertだけがDBに問い合わせ，観測行を返す．Final synthesisは，複数の近傍行が存在する場合，親子関係や時刻順，対象オブジェクトの一致が証跡で支えられるかを確認して主行動列を構成する．

\subsection{出力形式}

表\ref{tab:output}に，出力フィールドと評価上の扱いを示す．行動ステップ列（実装上の\code{code\_steps}）は主出力であり，採点時には個別の行動主張へ分解される．短縮行動系列（実装上の\code{code\_sequence}）は，観測されたコマンドまたは対象操作だけを短く並べた系列であり，アラート名やreasonの言い換えを含めない．\code{evidence}は行動を支えるログ証跡である．\code{excluded\_nearby\_evidence}は，時刻や対象が近いが主行動列には含めない行を明示するために用いる．

\begin{table}[t]
\centering
\small
\caption{出力フィールドと評価上の扱い}
\label{tab:output}
\begin{tabular}{p{0.30\linewidth}p{0.60\linewidth}}
\toprule
Field & 役割 \\
\midrule
\code{steps} & 主体，操作，対象，コマンドライン，証跡を持つ行動ステップ \\
\code{sequence} & 観測された操作を時系列に短く表した系列 \\
\code{evidence} & source stream，row id，timestamp，PID，command lineなどの根拠 \\
\code{nearby} & 近傍にあるが主系列ではない証跡 \\
\code{limitations} & 観測できない値，未確認の関係，復元上の限界 \\
\bottomrule
\end{tabular}
\end{table}

\section{実験設定}

\subsection{対象データセット}

実験には，ATLASv2 benign h1のWindowsホストログをCLOUSEAU用に変換したSQLiteデータベースを用いた．データベースには，正規化ログを格納する\code{audit\_logs}，DNS要求を格納する\code{dns\_requests}，ブラウザ履歴を格納する\code{browser\_history}が含まれる．\code{audit\_logs}には，Security，Sysmon，CBC EDR/NGAV events，CBC alertsを正規化して格納し，\code{source\_stream}により\code{msft-security}，\code{sysmon}，\code{cbc-edr}，\code{cbc-ngav}，\code{cbc-edr-alerts}，\code{cbc-ngav-alerts}などの出所を区別する．

プレ実験で用いた変換済みDBは，\code{audit\_logs} 7,531,007行，うちSysmon由来6,406行，CBC event telemetry由来1,010,918行，CBC alert summary由来232行を含む．また，\code{dns\_requests}は82,410行，\code{browser\_history}は53,069行である．DBの主要カラムは，time，pid，ppid，process name，access，object，event id，source stream，command lineであり，証跡保持のためにparent process，parent command line，process GUID，alert name，alert reason，registry/file/network関連フィールド，source row idも保持する．

\subsection{対象事例}

プレ実験では，Discord Run Key persistence事例を対象とした．調査起点は\code{host=WIN-32-H1}，\code{process=reg.exe}，\code{timestamp=2022-07-16 15:07:46}である．正解行動列は3ステップからなる．第一に，\code{Discord.exe}が\code{reg.exe}を子プロセスとして起動する．第二に，\code{reg.exe query}がユーザRun Key配下の\code{Discord}値を照会する．第三に，\code{reg.exe add}が同じ値へ\code{Update.exe --processStart Discord.exe}を指すデータを追加または更新する．

この事例は，Run Keyという明確な対象，親子プロセス，queryとaddの操作，レジストリオブジェクト，コマンドラインが評価可能であるため，プレ実験として扱いやすい．ただし，単一host，単一chain，単一attack patternの結果であるため，この結果だけから一般的なエンドポイント調査性能を主張することはできない．

\subsection{比較条件}

表\ref{tab:stages}に，3つの比較条件を示す．Stage 1は，CBC alert triage情報とhost/process/timestampを初期入力に含める実運用寄りの条件である．Stage 2は，初期入力からアラート情報を外し，host/process/timestampのみを与えるが，DB内にはCBC alert summaryとtelemetryを残す．Stage 3は，初期入力をStage 2と同じにしつつ，DBからCBC alert summary rowsのみを除外する．ただし，CBC EDR/NGAV event telemetry，Security，Sysmon，DNS，browser historyは残す．

\begin{table*}[t]
\centering
\small
\caption{プレ実験における3条件}
\label{tab:stages}
\begin{tabular}{p{0.12\linewidth}p{0.25\linewidth}p{0.35\linewidth}p{0.18\linewidth}}
\toprule
Stage & 初期入力 & DB条件 & 確認したい点 \\
\midrule
Stage 1 & CBC alert triage + host/process/timestamp & alert summary，CBC telemetry，OS側ログを保持 & アラート起点の実運用ベースライン \\
Stage 2 & host/process/timestampのみ & Stage 1と同じ全ログ保持DB & 少数clueから関連証跡へ到達できるか \\
Stage 3 & host/process/timestampのみ & CBC alert summary rowsのみ除外し，telemetryとOS側ログは保持 & アラート要約なしで統合できるか \\
\bottomrule
\end{tabular}
\end{table*}

Stage 3でalert summary rowsを外す理由は，モデルを不自然に不利にするためではない．CBC alert summaryには，alert name，reason，対象プロセス，対象レジストリ，検知理由など，調査者にとって答えに近い情報が集約されている場合がある．この状態だけを評価すると，モデルが低レベルログを探索・統合できたのか，alert summaryを言い換えただけなのかを区別しにくい．そこでStage 3では，高密度な要約行だけを外し，低レベルtelemetryからの探索能力を見る．

\subsection{実行設定}

比較対象モデルは，\code{gpt-4.1-mini}と\code{gpt-5.4-mini}の2種類であり，実験時点のAPI識別子をそのまま記録した．各モデルについて，Stage 1からStage 3までを1回ずつ実行し，合計6 runを得た．Stage 1はalert入力条件，Stage 2はprocess-time入力条件，Stage 3はprocess-time入力条件に加えてCBC alert summary除外条件として実行した．各runでは最大調査数100，最大質問数200，最大SQL問い合わせ数400，最大出力トークン8192を上限とした．Stage 3では，物理DBを削除するのではなく，SQL toolが参照するviewからCBC alert summary rowsを隠すことで条件を作った．このとき232行のalert summaryは検索対象から外れ，CBC event telemetry 1,010,918行は利用可能なまま残した．

各runの出力は，モデル名，入力条件，run JSON，採点結果を対応付けて保存した．本稿では，2026年6月6日に検証済みの6 runを対象とし，集計値ではなく各条件の結果をそのまま報告する．なお，本実験の目的はモデルの一般性能比較ではなく，入力情報とDB内要約行の有無が行動列復元に与える影響を予備的に確認することである．

\subsection{評価方法}

評価には，行動主張単位の採点指標を用いる．Goldは，対象事例のログ証跡を人手で確認し，process creation，registry query，registry add/updateの3ステップに正規化して作成した．各ステップについて，ログから直接確認できる主体，操作，対象，コマンドライン，critical evidenceを必須項目に含めた．推定されたユーザ意図，悪性判定，ファイル内容，アプリケーションの業務目的など，観測行だけから直接主張できない項目は必須項目に含めない．

本プレ実験では，正解行動列を合計14個の必須項目として固定した．D1は\code{Discord.exe}が\code{reg.exe}を子プロセスとして起動する行動であり，4項目を持つ．D2は\code{reg.exe query}によるRun Keyの\code{Discord}値照会であり，5項目を持つ．D3は\code{reg.exe add}によるRun Keyの\code{Discord}値追加または更新であり，5項目を持つ．主指標のrecallとprecisionはこの14項目を分母とする．また，候補側精度は候補出力に現れた必須項目相当の主張のうちgoldに対応した割合を表し，overclaim slot countはgoldに対応しない過剰な主張数を表す．

採点では，候補出力に含まれる主体，操作，対象，コマンドライン，critical evidenceを正規化し，goldの各項目と照合した．同じ意味を持つ表記揺れは許容するが，alert nameやalert reasonを操作そのものとして出力した場合は正解としない．証跡は，該当するログ行または同一行に含まれるmaterial termsで支えられる場合にのみ正解とした．過剰主張は，gold chainに属さない近傍行動，同一行動の重複表現，証跡で支えられない推定，および不完全なcommand lineやregistry pathの一般化として数えた．

\section{結果}

表\ref{tab:result}に，プレ実験の結果を示す．数値は，2026年6月6日に同一の採点規則で再採点し，6 runすべての入力・出力・採点ファイルの対応を確認した結果に基づく．本稿では，単一事例の結果であることを明確にするため，複数runの平均値ではなく，各モデルと条件の実行結果をそのまま報告する．

\begin{table*}[t]
\centering
\scriptsize
\caption{Discord Run Key事例におけるプレ実験結果．claim prec.は候補側精度，crit. ev.はcritical evidenceを表す}
\label{tab:result}
\begin{tabular}{llllllll}
\toprule
model & stage & condition & recall & claim prec. & order & crit. ev. & overclaim \\
\midrule
gpt-4.1-mini & Stage 1 & CBC alert input & 4/14 & 4/19 & 0/2 & 1/3 & 15 \\
gpt-4.1-mini & Stage 2 & all logs retained & 0/14 & 0/10 & 0/2 & 0/3 & 10 \\
gpt-4.1-mini & Stage 3 & alert summary removed & 0/14 & 0/17 & 0/2 & 0/3 & 17 \\
gpt-5.4-mini & Stage 1 & CBC alert input & 13/14 & 13/26 & 1/2 & 2/3 & 13 \\
gpt-5.4-mini & Stage 2 & all logs retained & 14/14 & 14/32 & 1/2 & 3/3 & 18 \\
gpt-5.4-mini & Stage 3 & alert summary removed & 8/14 & 8/18 & 0/2 & 1/3 & 10 \\
\bottomrule
\end{tabular}
\end{table*}

gpt-5.4-miniは，Stage 1で13/14の必須項目を復元した．これは，alert triage情報が初期入力として与えられる場合に，対象行動の大半を抽出できたことを示す．ただし，行動順序は1/2であり，すべての行動要素が正しく拾われていても，時系列または因果的な構成が完全になるとは限らない．

Stage 2では，gpt-5.4-miniのrecallは14/14となった．この条件では，アラート情報を初期入力から外しているため，モデルはhost/process/timestampから探索を開始する．それでも，DB内にalert summaryと関連telemetryが残っている場合には，正解行動要素へ到達できた．一方で，候補側精度は14/32であり，overclaimは18であった．したがって，主行動要素の復元と，余分な近傍行動を抑制することは別の課題である．

Stage 3では，gpt-5.4-miniのrecallは8/14に低下し，行動順序は0/2となった．CBC alert summary rowsを除外しても，EDR telemetryやOS側ログは残っているため，少なくとも一部の必須項目は観測可能である．しかし，query stepやcritical evidenceの同定が崩れ，Discord Run Key周辺の近いレジストリ操作が混入した．この結果は，alert summaryが最終出力として直接使われていない場合でも，探索経路と証跡同定の足場として機能している可能性を示している．

gpt-4.1-miniは，Stage 1で4/14を復元したが，Stage 2およびStage 3ではgoldの必須項目に対応する主張を復元できなかった．今回の設定では，少数のprocess-time clueから関連ログを探索し，複数証跡を統合する能力において，gpt-5.4-miniとの差が大きく表れた．ただし，この比較は単一事例に限られるため，モデル系列一般の優劣としては扱わない．

\section{考察}

\subsection{アラート要約の役割}

プレ実験から得られる重要な示唆は，alert summary rowsの有無が対象行動への到達性に影響した可能性である．Stage 1ではアラート情報が初期入力に含まれ，Stage 2では初期入力には含まれないがDB内に残っている．gpt-5.4-miniはこの2条件で高いrecallを示した．一方，Stage 3でalert summary rowsだけを除外すると，低レベルtelemetryが残っていても性能が低下した．

この結果は，アラート要約が「答え」そのものとして使われたことを意味しない．本手法では，アラート名やreasonを最終行動として出力することを避け，プロセス，コマンドライン，レジストリ対象，source rowに基づく行動を要求している．むしろ，alert summaryは，探索対象の時刻範囲や関連プロセス，関連レジストリを絞るためのインデックスとして機能したと解釈できる．今後は，Stage 2とStage 3のSQL履歴を比較し，どの段階で対象rowへ到達できなくなったかを分離して分析する必要がある．

\subsection{過剰主張と行動順序}

Stage 2のgpt-5.4-miniは，必須項目をすべて満たした一方で，候補側精度は14/32であった．これは，探索型エージェントが関連しそうなログ行を広く拾うほど，正解要素だけでなく近傍行動や重複主張も出力しやすいことを示す．SOC支援システムとしては，候補を広く提示することに価値がある場合もあるが，論文上の行動列復元としては，主行動列と補助証跡を明確に分ける必要がある．

また，Stage 2ではrecallが14/14であっても，orderは1/2に留まった．行動要素を個別に拾うことと，D1からD3までの順序関係を正しく構成することは異なる．特に，queryとaddが近い時間帯に発生し，同じプロセスやレジストリ対象を共有する場合，モデルは要素を列挙できても系列としての構造を誤る可能性がある．今後は，短縮行動系列だけでなく，行動ステップ列内のtimestamp，parent-child relation，evidence orderingを用いた順序制約を強化する．

\subsection{本実験への拡張}

本稿の結果はプレ実験であり，単一ケースに限定される．正式実験用の評価集合として，27個のbehavior chain，Stage 1/2用の75 behavior stepsおよび300必須項目，Stage 3用の65 answerable behavior stepsおよび260必須項目を準備している．各chainのgoldは，起点となる観測行の特定，process actionまたはobject operationへのstep分解，必須項目定義，critical evidence slotの固定，gold chainに属さない近傍行動の除外という手順で作成する．本稿では，この評価集合全体の結果ではなく，評価手順を単一事例で確認した結果のみを扱う．

この拡張により，Discord Run Key固有の観察にとどまらず，DNS packet capture，Python simple HTTP server，Sublime Text経由のscript execution，cmd.exe周辺操作など，複数の行動パターンで同じ評価枠組みが機能するかを検証できる．Stage 3の低下が複数ケースで再現する場合には，alert summaryに依存しない低レベルtelemetry探索，主行動列と近傍行動の分離，および順序構成の改善が中心課題となる．

\section{おわりに}

本稿では，少数のSOC調査起点からWindowsエンドポイントログを探索し，証跡付き行動列を復元するエージェント型ログ探索手法の評価枠組みを設計し，単一事例で予備的に検証した．Discord Run Key persistence事例のプレ実験では，gpt-5.4-miniが全ログ保持条件で14/14の必須項目を復元した一方，行動順序は1/2に留まり，過剰主張も残った．また，CBC alert summary rowsを除外した条件では8/14へ低下し，alert summaryが探索経路と証跡同定に寄与する可能性が示された．

今後は，27 chainの正式実験を完了し，複数host，複数process chain，複数行動パターンにおいて，Stage 1/2/3の差を定量的に分析する．あわせて，低レベルtelemetryのみからのquery planning，主行動列と近傍行動の分離，および時系列順序制約の改善を進める．

\begin{thebibliography}{9}
\bibitem{mitre_t1547_001}
MITRE ATT\&CK,
``Boot or Logon Autostart Execution: Registry Run Keys / Startup Folder, T1547.001,''
\url{https://attack.mitre.org/techniques/T1547/001/}
（参照 2026-06-09）．

\bibitem{sysmon}
Microsoft,
``Sysmon - Sysinternals,''
\url{https://learn.microsoft.com/en-us/sysinternals/downloads/sysmon}
（参照 2026-06-09）．

\bibitem{react}
S. Yao, J. Zhao, D. Yu, N. Du, I. Shafran, K. Narasimhan, and Y. Cao,
``ReAct: Synergizing Reasoning and Acting in Language Models,''
arXiv:2210.03629, 2022.

\bibitem{deeplog}
M. Du, F. Li, G. Zheng, and V. Srikumar,
``DeepLog: Anomaly Detection and Diagnosis from System Logs through Deep Learning,''
Proc. ACM CCS, pp. 1285--1298, 2017.

\bibitem{atlas}
A. Alsaheel, Y. Nan, S. Ma, L. Yu, G. Walkup, Z. B. Celik, X. Zhang, and D. Xu,
``ATLAS: A Sequence-based Learning Approach for Attack Investigation,''
Proc. USENIX Security, pp. 3005--3022, 2021.

\bibitem{atlasv2}
A. Riddle, K. Westfall, and A. Bates,
``ATLASv2: ATLAS Attack Engagements, Version 2,''
arXiv:2401.01341, 2023.
\end{thebibliography}

\end{CJK}
\end{document}
